\documentclass[conference]{IEEEtran}

\usepackage{iftex}
\ifPDFTeX
\usepackage[utf8]{inputenc}
\usepackage[T1]{fontenc}
\fi
\usepackage{amsmath,amssymb,bm,mathtools}
\usepackage{newtxtext,newtxmath}
\usepackage{textcomp}
\usepackage{array,booktabs}
\usepackage{graphicx}
\usepackage{siunitx}
\usepackage{cite}
\usepackage[breaklinks=true,hidelinks]{hyperref}
\usepackage{bookmark}
\usepackage[nameinlink,capitalise,noabbrev]{cleveref}
\usepackage[final]{microtype}

\interdisplaylinepenalty=2500
\allowdisplaybreaks
\emergencystretch=2em
\sisetup{detect-weight=true,detect-family=true,per-mode=symbol}

\DeclareRobustCommand{\vct}[1]{\bm{#1}}
\DeclareRobustCommand{\mat}[1]{\bm{#1}}
\newcommand{\T}{^{\mathsf T}}
\newcommand{\R}{\mathbb{R}}
\newcommand{\norm}[1]{\left\lVert #1\right\rVert}
\newcommand{\abs}[1]{\left|#1\right|}
\DeclareMathOperator{\diag}{diag}
\DeclareMathOperator{\wrap}{wrap}
\DeclareMathOperator{\sat}{sat}

\title{Separation of Steady Wind Heel and Wave-Induced Roll Using Marine IMU Sensor Fusion\\
\large State Decomposition, Apparent-Wind Aiding, Observability, and a Validation Program (concept paper)}
\author{%
  \IEEEauthorblockN{Mikhail S. Grushinskiy}
  \IEEEauthorblockA{Independent Researcher, USA, Fair Lawn\\
  Bareboat Necessities GitHub Project\\
  Email: mgrouch@users.github.com}%
}

\begin{document}
\maketitle

\begin{abstract}
A sailing vessel can carry a substantial quasi-steady roll angle under aerodynamic
load while simultaneously executing wave-induced roll about that moving equilibrium.
A conventional attitude estimator observes the physical attitude but does not, by
itself, identify which part of roll is steady wind heel and which part is wave motion.
That distinction matters when an IMU is used not only for navigation but also for wave
measurement: subtracting too much roll suppresses the signal of interest, while leaving
steady heel inside the wave frame distorts gravity compensation, horizontal acceleration,
and directional-motion estimates. This concept paper formulates the separation as a
latent-state sensor-fusion problem. The physical body frame is decomposed into a slowly
varying heel-neutral frame and a dynamic wave-attitude frame. The mean-heel state is
modeled as a slow response to an apparent-wind heeling-moment proxy, while a quaternion
attitude estimator retains the dynamic roll, pitch, and yaw response. Apparent wind is
used only as probabilistic low-frequency aiding because wind speed and angle do not
uniquely determine heel without sail trim, aerodynamic coefficients, hydrostatic
stiffness, and vessel speed. The paper identifies a central observability ambiguity:
low-frequency transverse acceleration and roll-induced gravity projection enter the
accelerometer through nearly the same channel, so total attitude sensing does not by
itself make the heel/wave decomposition observable. Gyroscope dynamics, magnetic
heading, time-scale separation, wind forcing, maneuver logic, and optional external
velocity/position aiding must supply the missing information. No performance result is
claimed in this draft. Instead, a predeclared program of constant-wind, gust, tack/gybe,
and irregular-directional-sea experiments is proposed, together with baselines, metrics,
and failure tests needed before the decomposition can be treated as a validated marine
estimation method.
\end{abstract}

\begin{IEEEkeywords}
sailboat heel, wave-induced roll, marine IMU, sensor fusion, apparent wind, attitude
estimation, Kalman filtering, observability, maneuver detection, wave measurement
\end{IEEEkeywords}

\section{Introduction}
\label{sec:introduction}

Marine IMUs are often asked to serve two purposes at once: estimate vessel attitude for
navigation and control, and estimate vessel motion as a proxy for the surrounding wave
field. The two purposes become coupled on a sailing vessel. Aerodynamic side force on
the sails produces a persistent heeling moment, so the mean hull roll can remain many
degrees away from zero even when the wave-induced roll has zero mean. A wave-motion
estimator that interprets the full roll angle as wave motion can contaminate gravity
compensation and horizontal acceleration. A preprocessing step that simply subtracts a
low-pass roll estimate can create the opposite error by removing genuine long-period
wave roll.

The repository already contains an optional implementation mechanism for an externally
supplied steady heel angle: body-frame accelerometer, gyroscope, and magnetometer vectors
can be transformed into a virtual un-heeled frame before the attitude and motion updates.
That mechanism has not been validated in the reported OU-filter experiments, and no
existing result establishes how the heel angle itself should be estimated. The present
paper is therefore deliberately prospective. It asks what estimator structure and what
evidence would be required to turn that frame transform into a defensible separation of
steady wind heel and dynamic wave attitude.

Attitude estimation from inertial and magnetic measurements is well established
\cite{Markley2003AttitudeError,MarkleyCrassidis2014_Attitude,Sola2017_QuaternionESKF},
and vessel-motion estimation from strapdown inertial sensors has been studied for heave
and directional wave applications \cite{Auestad2013_HeaveStrapdownIMU,Pascoal2009_KalmanDirectionalSpectrum,YurovskyKudinov2025_IMUwaves}.
The problem considered here is narrower and different: the physical attitude may be
well estimated while its decomposition into a slowly varying aerodynamic equilibrium and
a wave-induced perturbation remains ambiguous. That decomposition is not a new attitude
reference; it is a model-based partition of one measured motion into two latent causes.

The proposed contribution has three layers. First, the frame algebra is made explicit so
that changes in the heel estimate do not create artificial jumps in physical attitude.
Second, a slow heel state, a dynamic wave-attitude state, and an apparent-wind forcing
model are combined with explicit observability caveats. Third, the paper predeclares a
validation program broad enough to test the difficult cases: low-frequency horizontal
acceleration, gusts whose time scale overlaps the wave band, tack and gybe transitions,
and irregular directional seas. The objective is not to prove that frequency separation
always works; it is to identify when the separation is identifiable, when it is merely a
useful prior, and when the estimator must report ambiguity rather than invent a clean
heel/wave split.

\section{Sailboat Heel and Wave-Roll Dynamics}
\label{sec:dynamics}

Let the physical roll of the hull be generated by at least two mechanisms. The first is
a slowly varying equilibrium under aerodynamic heeling moment and hydrostatic restoring
moment. The second is a dynamic response to waves, vessel maneuvering, gust transients,
and coupled roll--yaw--sway dynamics. For modest pitch and local analysis it is tempting
to write
\begin{equation}
 \phi_{\mathrm{body}}(t)\approx \phi_h(t)+\phi_w(t),
 \label{eq:additive-roll}
\end{equation}
where $\phi_h$ is steady or slowly varying heel and $\phi_w$ is wave-induced roll. The
matrix formulation used later is preferable because rotations do not add globally, but
\eqref{eq:additive-roll} captures the intended time-scale decomposition.

A simple quasi-static roll balance is
\begin{equation}
 M_{\mathrm{aero}}(t)+M_{\mathrm{wave,LF}}(t)
 \approx M_{\mathrm{rest}}\!\left(\phi_h(t)\right),
 \label{eq:heel-balance}
\end{equation}
where the low-frequency wave term is written explicitly to emphasize an important
limitation: a long swell or persistent asymmetric wave loading can shift the apparent
mean roll even without wind. Near an operating point, the hydrostatic restoring moment
may be linearized as $M_{\mathrm{rest}}\approx K_\phi\phi_h$, but $K_\phi$ depends on
loading, geometry, and heel itself.

The dynamic component can be represented at several levels of fidelity. A minimal
single-mode roll model is
\begin{equation}
 \ddot\phi_w+2\zeta_r\omega_r\dot\phi_w+\omega_r^2\phi_w
 = b_r u_{\mathrm{wave}} + w_r,
 \label{eq:roll-oscillator}
\end{equation}
where $\omega_r$ and $\zeta_r$ are effective roll parameters. For a general directional
sea this scalar oscillator is only a shaping prior: the actual vessel response contains
multiple encounter frequencies and coupling to pitch, yaw, sway, and surge. A full
quaternion attitude state is therefore retained in the proposed sensor-fusion
architecture; \eqref{eq:roll-oscillator} is useful as a hypothesis to test, not as a
required truth model.

The desired separation is consequently asymmetric. The steady-heel state should be
allowed to move when the aerodynamic equilibrium really moves, but it should not chase
ordinary wave roll. The dynamic attitude state should follow rapid roll accurately, but
it should not be forced to carry a large quasi-static heel if the downstream wave-motion
model assumes an approximately zero-mean roll frame.

\section{Coordinate Frames and Measurements}
\label{sec:frames}

Use a world NED frame $\mathcal W$, a physical IMU/body frame $\mathcal B$, and a virtual
heel-neutral frame $\mathcal B'$. The body axes are $x$ forward, $y$ starboard, and $z$
down. Positive roll follows the right-hand rule about $+x$. Let
$\mat R_{WB}$ map vector coordinates from $\mathcal B$ to $\mathcal W$ and let
$\mat R_{WB'}$ map $\mathcal B'$ to $\mathcal W$. Define the mean-heel transform by
\begin{equation}
 \vct v_{B'}=\mat R_x(-\phi_h)\vct v_B.
 \label{eq:deheel-vector}
\end{equation}
Then the physical attitude is
\begin{equation}
 \boxed{\mat R_{WB}=\mat R_{WB'}\mat R_x(-\phi_h).}
 \label{eq:attitude-compose}
\end{equation}
The dynamic wave roll is retained inside $\mat R_{WB'}$; only the slowly varying
$\phi_h$ is removed from the physical sensor frame.

The IMU measurements are modeled conventionally as
\begin{align}
 \vct\omega_m &= \vct\omega_B+\vct b_g+\vct n_g,\\
 \vct f_m &= \mat R_{BW}(\vct a_W-\vct g_W)+\vct b_a+\vct n_a,\\
 \vct m_m &= \mat R_{BW}\vct m_W+\vct b_m+\vct n_m,
 \label{eq:imu-model}
\end{align}
with the usual caution that magnetic disturbances can be strongly vessel dependent. The
accelerometer is a specific-force sensor, not a gravity-direction sensor whenever vessel
translation is appreciable. This distinction becomes central in
\cref{sec:observability}.

The apparent-wind instrument supplies at least a relative wind angle $\beta_a$ and often
a speed $V_a$. If a vector measurement is available at a masthead lever arm $\vct r_m$,
then rotational velocity contributes to the local apparent wind. A first-order correction
has the structure
\begin{equation}
 \vct v_{a,m}
 \simeq \vct v_{a,O}-\vct\omega_B\times\vct r_m+\vct n_w,
 \label{eq:wind-lever}
\end{equation}
with sign depending on the chosen relative-velocity convention. This term should be
modeled or its uncertainty inflated during large roll/yaw rates; otherwise the wave
motion being estimated can feed back into the wind-based heel aid.

\section{Steady-Heel State/Model}
\label{sec:heel-model}

A useful first model is deliberately low order. Let
\begin{equation}
 \vct x_h=\begin{bmatrix}\phi_h & \kappa_h\end{bmatrix}\T,
 \label{eq:heel-state}
\end{equation}
where $\kappa_h$ is an effective mapping from a wind-load proxy to equilibrium heel. The
state obeys
\begin{align}
 \dot\phi_h &= -\frac{1}{\tau_h}\left(\phi_h-\phi_{\mathrm{eq}}\right)+w_h,\\
 \dot\kappa_h &= w_\kappa,
 \label{eq:heel-dynamics}
\end{align}
with a long heel time constant $\tau_h$ relative to the dominant wave period under
steady conditions. A candidate equilibrium map is
\begin{equation}
 \phi_{\mathrm{eq}}
 = \sat_{\phi_{\max}}\!\left(\kappa_h u_w\right),
 \qquad
 u_w=V_a^2 g_w(\beta_a,\delta_s,V_b),
 \label{eq:heel-equilibrium}
\end{equation}
where $\delta_s$ denotes sail trim if available and $V_b$ vessel speed. The exact function
$g_w$ is not prescribed in this draft. A signed lateral-force proxy such as
$\sin\beta_a$ is a useful starting hypothesis, but it must be identified against vessel
and sail configuration rather than treated as universal aerodynamics.

An alternative second-order state
\begin{equation}
 \ddot\phi_h+2\zeta_h\omega_h\dot\phi_h
 +\omega_h^2(\phi_h-\phi_{\mathrm{eq}})=w_h
 \label{eq:heel-second-order}
\end{equation}
may be required when gust-driven heel transients are not slow compared with wave roll.
The experiments in \cref{sec:gusts} should determine whether the added state improves
separation or merely gives the heel model enough bandwidth to absorb wave energy.

The process covariance of $\phi_h$ is therefore a scientific parameter, not just a
numerical tuning constant. Too little process noise prevents adaptation to real wind
changes; too much makes the ``steady'' heel state a low-pass copy of the total roll.

\section{Dynamic Wave-Attitude Model}
\label{sec:wave-attitude}

The dynamic attitude is represented by a quaternion or rotation matrix for
$\mathcal B'\rightarrow\mathcal W$ and the usual local attitude-error state
\cite{Markley2003AttitudeError,Sola2017_QuaternionESKF}. A minimal error-state vector is
\begin{equation}
 \delta\vct x_a=
 \begin{bmatrix}
   \delta\vct\theta & \vct b_g & \vct b_a
 \end{bmatrix}\T,
 \label{eq:att-state}
\end{equation}
optionally augmented by translational wave states. The gyroscope update is performed on
the de-heeled angular-rate vector, while accelerometer and magnetometer predictions are
formed in $\mathcal B'$. The physical hull attitude returned to applications is restored
through \eqref{eq:attitude-compose}.

The crucial design rule is that the attitude estimator must not be asked to zero the
dynamic roll. The wave component is a legitimate part of $\mat R_{WB'}$. In a
frequency-aware implementation, the existing wave-period estimator can provide a slowly
varying characteristic frequency $\omega_p$ that shapes, but does not dictate, the
separation prior. One candidate is to model the wave-roll innovation with a band-limited
Gauss--Markov or second-order process centered near encounter frequency while leaving the
quaternion kinematics exact.

The dynamic model should also preserve cross-axis motion. Directional waves can generate
roll, pitch, yaw, and horizontal acceleration together; forcing an independent scalar
roll oscillator can improve one metric while breaking the cross-covariances that a
three-dimensional inertial estimator needs. Accordingly, the scalar roll model in
\eqref{eq:roll-oscillator} is best treated as an ablation against the full 3-D attitude
model.

\section{Apparent-Wind Aiding}
\label{sec:wind-aiding}

Apparent wind is informative about heel because aerodynamic side force is a major cause
of steady sailing heel, but it is not a direct angle measurement. A schematic heeling
moment is
\begin{equation}
 M_a \approx \frac12\rho_{air}V_a^2 A_s h_{CE}
 C_H(\beta_a,\delta_s,\mathrm{reef},\ldots),
 \label{eq:wind-moment}
\end{equation}
so the same $V_a$ and $\beta_a$ can produce different heel for different sail trim,
reefing, loading, righting moment, and boat speed. The proposed filter therefore uses
wind as a broad, adaptive prior on $\phi_h$, not as a hard pseudo-measurement.

One formulation writes
\begin{equation}
 z_w=0=\phi_h-\kappa_h u_w+\nu_w,
 \qquad \nu_w\sim\mathcal N(0,R_w),
 \label{eq:wind-pseudo}
\end{equation}
with a deliberately large $R_w$ and slowly adapting $\kappa_h$. A more conservative
form uses only the expected sign of heel and a magnitude envelope. Both should be tested.
A hard mapping from apparent wind to heel is unlikely to be robust enough across reef
changes, crew movement, sail trim, and different boats.

Three additional issues deserve explicit experiments. First, the wind sensor is mounted
on the moving vessel and can contain wave-induced rotational-velocity contamination as in
\eqref{eq:wind-lever}. Second, apparent wind changes when vessel speed and yaw change even
if true wind does not. Third, sensor latency can be comparable with a gust transient.
Wind aiding should therefore carry time alignment, lever-arm correction when possible,
and dropout/innovation gating.

\section{Heel-Frame Transformation}
\label{sec:heel-transform}

The frame transform is the mechanism that allows the downstream wave-motion estimator to
operate in a heel-neutral body frame without deleting wave roll. Every body-frame vector
used by that estimator is mapped as in \eqref{eq:deheel-vector}, while the dynamic
attitude remains $\mat R_{WB'}$.

Changing $\hat\phi_h$ changes the definition of $\mathcal B'$. If the heel estimate is
updated by $\Delta\phi_h$, the internal attitude must be retargeted so that the physical
attitude does not jump. From \eqref{eq:attitude-compose}, invariance of $\mat R_{WB}$
requires
\begin{equation}
 \boxed{
 \mat R_{WB'}^{+}=\mat R_{WB'}^{-}\mat R_x(\Delta\phi_h).}
 \label{eq:retarget}
\end{equation}
The covariance must be transformed consistently with the local attitude coordinates.
This is especially important if $\phi_h$ changes rapidly after a tack: merely changing
the preprocessing rotation while leaving the quaternion untouched creates an artificial
attitude step.

The repository's current optional heel path already follows this conceptual pattern: an
externally commanded heel angle retargets the virtual body frame before subsequent IMU
updates. The missing problem is estimation and validation of that command. A first unit
test for the new estimator should therefore be purely geometric: arbitrary sequences of
$\phi_h$ updates with no physical motion must leave the recovered physical quaternion,
gravity prediction, and magnetic prediction unchanged to numerical precision.

\section{Observability and Ambiguity with Horizontal Acceleration}
\label{sec:observability}

This is the central theoretical problem. Linearizing the accelerometer around an attitude
estimate gives an attitude-error contribution of the form
\begin{equation}
 \delta\vct f
 \simeq -\big[\mat R_{BW}(\vct a_W-\vct g_W)\big]_\times
 \delta\vct\theta
 +\mat R_{BW}\delta\vct a_W+\delta\vct b_a.
 \label{eq:acc-linearized}
\end{equation}
For small roll error and weak translation, the lateral channel contains approximately
$g\,\delta\phi$. The same channel also contains real transverse acceleration and
accelerometer bias. Therefore a slowly varying lateral specific-force residual can be
explained by changing mean heel, changing horizontal acceleration, or changing bias.
Those explanations are locally close to collinear in the measurement Jacobian.

A magnetometer and a gyroscope improve physical-attitude estimation, but they do not
fully solve the decomposition problem. The magnetometer observes total orientation with
respect to the magnetic field, not the semantic labels ``steady heel'' and ``wave roll.''
A constant heel has zero angular rate, so the gyro cannot determine its absolute offset.
The decomposition becomes identifiable only through additional structure: different
process time scales, a wind-correlated slow state, bounded bias dynamics, vessel-motion
models, and potentially GNSS or water-speed aiding.

The validation program should therefore include an explicit local observability study.
For an augmented state containing at least $(\delta\phi_w,\phi_h,b_{a,y},a_y)$, construct
the finite-horizon linearized observability Gramian
\begin{equation}
 \mat W_o(t_0,T)=
 \int_{t_0}^{t_0+T}
 \mat\Phi(t,t_0)\T\mat H(t)\T\mat R^{-1}(t)
 \mat H(t)\mat\Phi(t,t_0)\,dt,
 \label{eq:gramian}
\end{equation}
and report its smallest singular value and condition number across the study cases. The
same experiment should be repeated with wind aiding removed, with magnetometer removed,
and with horizontal GNSS velocity added. The purpose is not to obtain one global
observability theorem for a nonlinear sailing vessel; it is to identify which excitation
and aiding combinations actually separate the slow heel state from low-frequency
horizontal inertial error.

A particularly important negative control is constant transverse acceleration with no
wind change. An estimator that responds by changing $\phi_h$ has confused kinematics
with heel. Conversely, a constant wind-induced heel on calm water should not be absorbed
into accelerometer bias. These two cases are almost mirror images and should be treated as
primary identifiability tests.

\section{Adaptive Separation of Mean Heel and Wave Roll}
\label{sec:adaptive-separation}

The simplest adaptive architecture is a slow heel process coupled to a full dynamic
attitude filter. Let $\omega_p$ be a robust estimate of the dominant wave/encounter
frequency. A separation parameter can be expressed nondimensionally as
\begin{equation}
 \lambda_h=\frac{1/\tau_h}{\omega_p},
 \label{eq:lambda-h}
\end{equation}
with $\lambda_h\ll1$ under stationary sailing. Rather than fixing one cutoff in seconds,
the estimator can adapt $\tau_h$ so that the heel process remains slow relative to the
observed sea state. This is analogous to other sea-state-normalized schedules in the
repository, but here the design objective is spectral separation rather than integral
regularization.

A candidate update hierarchy is:
\begin{enumerate}
 \item estimate total physical attitude from gyro, accelerometer, and magnetometer;
 \item propagate a slow heel state using the wind-forced model;
 \item retarget $\mathcal B'$ continuously using \eqref{eq:retarget};
 \item keep wave roll in $\mat R_{WB'}$ rather than subtracting instantaneous physical roll;
 \item adapt $\tau_h$, $Q_h$, or $R_w$ only from slow statistics, with explicit bounds;
 \item suspend or widen the slow-state model during maneuvers and large innovations.
\end{enumerate}

Several variants should be compared rather than assumed equivalent: a fixed low-pass
heel estimate, a first-order Kalman heel state, the wind-aided state
\eqref{eq:wind-pseudo}, the second-order model \eqref{eq:heel-second-order}, and a joint
augmented-state EKF in which $\phi_h$ participates directly in the sensor Jacobians. The
joint formulation is theoretically cleaner but creates stronger coupling between heel,
attitude, horizontal acceleration, and bias. A cascade can be easier to stabilize and
interpret but may underestimate cross-covariance.

Adaptation should also report uncertainty. When the estimated slow and wave spectra
substantially overlap, a confidence measure can be based on the posterior variance of
$\phi_h$, the innovation consistency of the wind pseudo-measurement, and the observability
Gramian. The correct output in an ambiguous interval may be ``heel estimate uncertain,''
not an aggressively filtered angle.

\section{Maneuver/Tack Detection}
\label{sec:maneuvers}

Tacks and gybes deliberately violate the stationary-heel assumption. A tack typically
changes the sign of the aerodynamic heeling moment while yaw rate, apparent-wind angle,
and often roll rate all change together. A gybe can produce a large transient with a
different apparent-wind geometry. Treating either event as ordinary process noise can
make the slow heel state lag through the maneuver and temporarily contaminate the wave
frame.

A practical event detector should combine circular apparent-wind angle, yaw rate, and
heel evidence. One candidate state machine has stable-port, transition, and
stable-starboard states. A tack candidate is triggered when the signed apparent-wind
angle crosses the bow-sector centerline with sufficient yaw rotation; confirmation
requires dwell on the opposite side and, when wind loading is substantial, a compatible
change in the slow heel sign. A gybe uses an analogous circular crossing around the aft
sector rather than a naive discontinuity at $\pm180^\circ$.

During a confirmed maneuver the estimator can temporarily:
\begin{enumerate}
 \item increase $Q_h$ so the equilibrium heel can move;
 \item freeze adaptation of $\kappa_h$ so one maneuver does not relearn the aerodynamic map;
 \item inflate $R_w$ while the masthead wind is contaminated by rapid rotation;
 \item retain continuous physical attitude through the frame-retarget rule;
 \item restart the slow-separation statistics only after a post-maneuver dwell.
\end{enumerate}
The detector itself needs a scored study: event recall, false detections per hour,
detection latency, and the time for heel and wave-roll estimates to recover afterward.

\section{Validation Methodology}
\label{sec:validation}

No quantitative result is reported in this concept draft. The following validation plan
is intended to prevent post-hoc selection of only favorable sailing conditions.

The simulator should expose separate truth signals for $\phi_h(t)$ and the dynamic vessel
attitude, then compose them through \eqref{eq:attitude-compose} before synthesizing IMU,
magnetometer, and apparent-wind measurements. This is essential: if ``truth heel'' is
constructed afterward by low-pass filtering the same roll signal being evaluated, the
validation becomes circular.

At minimum the experiment should compare six estimators:
\begin{enumerate}
 \item no heel separation: the physical attitude is used directly;
 \item fixed-cutoff low-pass subtraction of roll;
 \item slow heel state without wind aiding;
 \item wind-aided adaptive heel state;
 \item an oracle heel transform supplied the true $\phi_h$;
 \item a deliberately over-fast heel state showing the failure mode in which wave roll is absorbed into heel.
\end{enumerate}
The oracle is not deployable; it establishes the maximum downstream benefit available if
mean heel were known exactly.

Each stochastic scenario should use paired seeds for wave phases, IMU noise/bias,
magnetometer error, and wind-sensor noise so that estimator differences are paired. A
multi-seed ensemble should be predeclared before examining results. The repository's
existing ten-seed paired-validation machinery provides a useful template, but the heel
study requires new scenario generators and new truth channels.

Primary metrics should include
\begin{align}
 e_h &= \hat\phi_h-\phi_h^{\mathrm{true}},\\
 e_w &= \hat\phi_w-\phi_w^{\mathrm{true}},
 \label{eq:errors}
\end{align}
reported as RMS, mean bias, 95th-percentile absolute error, and paired uncertainty across
seeds. Frequency-domain metrics are equally important. Define a wave-band leakage ratio
\begin{equation}
 L_{w\rightarrow h}=
 \frac{\int_{\Omega_w}S_{\hat\phi_h-\phi_h}(\omega)d\omega}
      {\int_{\Omega_w}S_{\phi_w}(\omega)d\omega},
 \label{eq:wave-to-heel-leak}
\end{equation}
and a low-frequency heel leakage into the wave estimate analogously. Report gain and
phase error of $\hat\phi_w$ over the truth wave-roll band, not only time-domain RMS.

The downstream effect should also be measured. The same physical run should be processed
through the wave-motion estimator with no heel correction, estimated heel, and oracle
heel. Compare horizontal/vertical acceleration residuals, displacement RMS, wave-height
statistics, and wave-direction metrics. A heel separator that improves its own scalar
RMSE but damages wave estimation has failed the actual application.

Finally, the validation must include sensor-model ablations: magnetometer dropout and
bias, apparent-wind latency, mast lever-arm error, accelerometer bias/random walk, heading
error, and optional GNSS velocity aiding. The observability claim should be conditioned on
which of those sensors are present.

\section{Constant-Wind Cases}
\label{sec:constant-wind}

The first study should remove maneuver and gust confounds. Use constant apparent-wind
forcing and impose equilibrium heel values spanning upright to strongly heeled sailing,
for example a grid in $\abs{\phi_h}$ from near zero through approximately
\SI{30}{\degree}. Both signs must be tested. Repeat each heel level with no waves, regular
single-frequency roll, and irregular seas of increasing energy.

The calm-water constant-heel cases answer the basic bias question: can the estimator
recover a static heel without converting it into accelerometer bias or dynamic roll? The
wave cases answer the leakage question: as the wave period approaches the chosen heel
time constant, how much true roll is absorbed by the slow state? A sweep over
$\lambda_h$ in \eqref{eq:lambda-h} should be reported rather than selecting one crossover
from a single sea state.

Wind-aiding ablations should vary the mapping error in $\kappa_h$, including an initially
wrong sign/magnitude and a mid-run sail-trim change with unchanged wind. These cases are
necessary to determine whether the wind model helps or merely imposes a plausible-looking
heel prior.

\section{Gusts and Changing Wind}
\label{sec:gusts}

The second study should vary wind forcing on time scales both slower and faster than the
wave band. Recommended inputs include smooth ramps, steps filtered by a realistic wind
sensor, isolated gust pulses, and stochastic colored gusts. The difficult regime is not
an extremely slow wind change; it is a gust with a characteristic period near the vessel
roll response or dominant wave period.

For every gust case report the peak heel error, settling time, maximum wave-roll
attenuation, and the fraction of gust-driven roll assigned to the dynamic state. There is
no universally correct semantic split during a fast gust: a gust produces real dynamic
roll even though its cause is wind. The paper should therefore distinguish
\emph{equilibrium heel} from \emph{wind-caused roll}. Only the former belongs in
$\phi_h$. A successful estimator should allow a rapid gust response to remain temporarily
in the dynamic attitude and then move the slow equilibrium as evidence accumulates.

A useful stress test is a wind step with no waves followed by the same wind step embedded
in an irregular sea. The difference measures whether wave excitation corrupts the slow
state adaptation. Another is a wind-direction change at nearly constant speed, which
tests the signed force proxy independently of the $V_a^2$ scaling.

\section{Tacks and Gybes}
\label{sec:tack-gybe-study}

The maneuver study should generate controlled tacks with known yaw trajectory, wind
crossing, and heel reversal. Separate slow tacks from aggressive tacks so that the heel
sign change occurs over a range of time scales. Gybes should be modeled separately
because the apparent-wind angle crosses the circular aft boundary and the roll transient
need not resemble a tack.

The primary estimator metric is recovery, not just event classification. Report tack/gybe
detection latency, false triggers, maximum physical-attitude discontinuity, peak error in
$\hat\phi_h$, peak error in wave roll, and time to recover the predeclared steady-state
error envelope. The geometric invariant in \eqref{eq:retarget} should be tested during
every event; physical attitude must remain continuous even if the estimated heel state is
allowed to change quickly.

Negative controls should include a large wave-roll excursion with no maneuver, a heading
change under power with little heel, and a short magnetometer disturbance. The detector
must not convert those events into tack resets merely because yaw or roll crosses a
threshold.

\section{Irregular Directional Seas}
\label{sec:directional-seas}

The final simulation stage should use irregular directional seas rather than only scalar
roll forcing. The existing JONSWAP and PM/Stokes generators provide a starting point for
broad sea-state amplitude and period sweeps, but the heel study needs vessel rotational
response and encounter geometry. At least the following should be included: beam seas,
head/following seas, quartering seas, and two-component crossing seas with different peak
periods and directions.

Directional seas matter for two reasons. First, true roll energy depends strongly on wave
direction and vessel response, so an adaptive separator tuned on beam-sea roll may become
overconfident in head seas. Second, horizontal acceleration and attitude are correlated;
this is precisely the regime in which the ambiguity of \cref{sec:observability} becomes
important. The study should therefore report not only roll error but also the horizontal
specific-force residual and the smallest singular value of the augmented observability
Gramian.

Long-period swell is a required edge case. If a swell component enters the nominal
``slow'' band, a purely spectral separator will classify real wave roll as heel. The
wind-aided model should improve that case only if the wind evidence is informative; if it
is not, uncertainty should increase. Crossing seas with one long-period component provide
an even stronger test of whether the method is truly model-based or merely a low-pass
filter with extra states.

\section{Results}
\label{sec:results}

\textbf{No experimental or simulation results are claimed in this draft.} This section
is intentionally a reporting contract for the studies above. It should not be populated
with deterministic examples and then described as validation.

The completed study should contain, at minimum, one primary table for constant-wind
steady-state separation, one for gust/change transients, one for tack/gybe recovery, and
one for directional-sea performance. Each table should report the paired estimator
comparison for heel RMS/bias, wave-roll RMS, wave-band gain/phase error, and both leakage
ratios. A separate table should report the observability diagnostics with and without
wind, magnetometer, and GNSS-velocity aiding.

Recommended figures are: (i) a time trace showing physical roll, truth mean heel, truth
wave roll, and both estimates; (ii) a PSD decomposition demonstrating leakage between the
slow and wave bands; (iii) a gust case with wind forcing and state uncertainty; (iv) a
tack showing frame retargeting without a physical-attitude jump; and (v) an observability
condition-number map versus wave direction, wind angle, and heel magnitude.

The primary claim should be accepted only if the adaptive method improves separation over
the fixed-cutoff baseline across the predeclared ensemble without materially degrading
total physical attitude or downstream wave-motion estimates. If improvement occurs only
when wind and wave spectra are widely separated, that narrower result should be stated
explicitly. If apparent-wind aiding helps only under some sail configurations, the method
should expose that dependency rather than absorb it into an opaque tuned coefficient.

\section{Discussion}
\label{sec:discussion}

The proposed architecture reframes ``heel correction'' as state decomposition. That
framing avoids a common conceptual error: the IMU does not observe two roll angles. It
observes one physical motion. Mean heel and wave roll become separate quantities only
after process assumptions and exogenous information are introduced. Consequently, good
physical-attitude accuracy is necessary but not sufficient evidence that the decomposition
is correct.

The most promising information source is not an accelerometer-derived static tilt angle;
that channel is precisely where low-frequency horizontal acceleration creates ambiguity.
Instead, the decomposition should combine the gyro's dynamic information, magnetometer or
other attitude reference, a deliberately slow heel process, and apparent wind as a noisy
causal input. External horizontal velocity or position aiding could be especially
valuable because it directly reduces the acceleration/tilt ambiguity identified in
\eqref{eq:acc-linearized}.

Apparent-wind aiding is also a possible failure source. Sail trim, reefing, crew movement,
and boat loading change the wind-to-heel map. A model that is too confident can make the
heel estimate look smooth while being systematically wrong. For this reason the wind
coefficient should adapt slowly, its covariance should widen after configuration changes,
and the estimator should retain a wind-free mode. If the method eventually requires sail
load sensors, traveler position, or hydraulic pressure to become reliable, those should
be treated as additional measurements rather than hidden tuning assumptions.

The virtual heel frame is useful even if the final estimator is not Kalman-based. The
geometric requirement \eqref{eq:retarget} applies equally to a complementary filter,
nonlinear observer, or Lie-group estimator: changing the definition of the heel-neutral
frame must not change the represented physical attitude. The same separation concept can
therefore be tested across OU-III, TFG, or other attitude/motion back ends once the
measurement and state semantics are fixed.

Real-vessel truth remains challenging. Simulation can define $\phi_h$ and $\phi_w$
separately by construction, but a field reference normally measures only total attitude.
A strong validation campaign should therefore combine several regimes: flat-water or
quasi-static sailing to establish mean-heel accuracy; a wave tank or controlled motion
platform where the decomposition is imposed; and open-water trials using a high-grade
attitude reference for total roll plus independent wind and navigation measurements. In
open water, any ``truth'' split produced only by post-filtering the reference roll should
be labeled an operational definition rather than ground truth.

\section{Conclusion}
\label{sec:conclusion}

Separating steady wind heel from wave-induced roll is not equivalent to estimating roll
accurately and then applying a low-pass filter. It is a latent-cause decomposition in
which the accelerometer contains a fundamental ambiguity between gravity projection,
horizontal acceleration, and bias. A defensible solution therefore needs explicit frame
algebra, a slow heel process, a dynamic three-dimensional attitude model, cautious
apparent-wind aiding, maneuver logic, and uncertainty that grows when the two components
cannot be distinguished.

The repository already has the geometric machinery required to work in a virtual
heel-neutral frame. The next step is evidence, not another tuning constant. The proposed
study program begins with constant-wind identifiability, then introduces gusts, tacks and
gybes, and finally irregular directional seas with sensor degradations and observability
analysis. The central questions are measurable: how much true wave roll leaks into the
heel state, how much slow heel contaminates the wave state, which aiding signals break
the horizontal-acceleration ambiguity, and whether the separation improves downstream
wave-motion estimation without degrading physical attitude. Until those questions are
answered over a paired multi-seed ensemble and controlled reference experiments, the
method should remain a structured hypothesis rather than a validated performance claim.

\bibliographystyle{IEEEtran}
\bibliography{w3d}

\end{document}
